\FigureLayoutDeclare{img/2010/hubei_science/q13_fig1.png}{3.0cm}{1bp 18bp 20bp 17bp}
\FigureLayoutDeclare{img/2010/hubei_science/q1_fig1.png}{8.0cm}{17bp 23bp 51bp 19bp}

\chapter{2010年湖北卷（理）}

\section{单选题}

\begin{problem}
若 \(\i\) 为虚数单位，图中复平面内点 \(Z\) 表示复数 \(z\)，则表示复数 \(\frac{z}{1+\i}\) 的点是（\quad）

\bitmapfigure[max width=0.50\linewidth]{img/2010/hubei_science/q1_fig1.png}

\choices
    {\(E\)}
    {\(F\)}
    {\(G\)}
    {\(H\)}
\end{problem}

\begin{answer}
D
\end{answer}

\begin{solution}
由图可读出点 \(Z\) 的坐标为 \((3,1)\)，所以 \(z=3+\i\).于是
\[
\frac{z}{1+\i}
=\frac{(3+\i)(1-\i)}{(1+\i)(1-\i)}
=\frac{4-2\i}{2}
=2-\i
\]
复数 \(2-\i\) 对应复平面内横坐标为 \(2\)、纵坐标为 \(-1\) 的点，即图中的 \(H\)，故选 D.
\end{solution}

\begin{problem}
设集合 \(A = \left\{(x, y) \mid \frac{x^2}{4} + \frac{y^2}{16} = 1\right\}\)，\(B = \{(x, y) \mid y = 3^x\}\)，则 \(A \cap B\) 的子集的个数是（\quad）

\choices
    {4}
    {3}
    {2}
    {1}
\end{problem}

\begin{answer}
A
\end{answer}

\begin{solution}
因为 \(3^x>0\)，交点只能位于椭圆的上半部分.令
\[
\phi(x)=\frac{x^2}{4}+\frac{9^x}{16}
\]
交点的横坐标正好是方程 \(\phi(x)=1\) 的根.又
\[
\phi''(x)=\frac12+\frac{(\ln9)^2 9^x}{16}>0
\]
所以 \(\phi\) 严格凸，方程 \(\phi(x)=1\) 至多有两个根.

另一方面，
\(\phi(-2)=1+\frac1{144}>1\)，\(\phi(0)=\frac1{16}<1\)，\(\phi(1)=\frac{13}{16}<1\).
当 \(x=\log_3 4\) 时，\(9^x=16\)，故
\[
\phi(\log_3 4)=\frac{(\log_3 4)^2}{4}+1>1
\]
由连续性，方程在 \((-2,0)\) 和 \((1,\log_3 4)\) 内各有一个根；结合至多两个根，恰有两个交点.因此 \(A\cap B\) 有两个元素，其子集个数为 \(2^2=4\)，故选 A.
\end{solution}

\begin{problem}
在 \(\triangle ABC\) 中，\(a = 15\)，\(b = 10\)，\(A = 60^{\circ}\)，则 \(\cos B =\)（\quad）

\choices
    {\(-\frac{2\sqrt{2}}{3}\)}
    {\( \frac{2\sqrt{2}}{3} \)}
    {\(-\frac{\sqrt{6}}{3}\)}
    {\(\frac{\sqrt{6}}{3}\)}
\end{problem}

\begin{answer}
D
\end{answer}

\begin{solution}
由正弦定理，
\(\frac{a}{\sin A}=\frac{b}{\sin B}\)，\(\sin B=\frac{b\sin A}{a} =\frac{10\sin60^\circ}{15} =\frac{\sqrt3}{3}\).
因 \(a>b\)，所以 \(A>B\)，从而 \(0<B<60^\circ\)，应取正的余弦值.因此
\[
\cos B=\sqrt{1-\sin^2B}
=\sqrt{1-\frac13}
=\frac{\sqrt6}{3}
\]
故选 D.
\end{solution}

\begin{problem}
投掷一枚均匀硬币和一枚均匀骰子各一次，记“硬币正面向上”为事件 \(A\)，“骰子向上的点数是 3”为事件 \(B\)，则事件 \(A\)，\(B\) 中至少有一件发生的概率是（\quad）

\choices
    {\(\frac{5}{12}\)}
    {\(\frac{1}{2}\)}
    {\(\frac{7}{12}\)}
    {\(\frac{3}{4}\)}
\end{problem}

\begin{answer}
C
\end{answer}

\begin{solution}
事件 \(A\) 与 \(B\) 相互独立，且
\(P(A)=\frac12\)，\(P(B)=\frac16\).
用对立事件，
\[
\begin{aligned}
P(A\cup B)
&=1-P(\overline A\cap\overline B)\\
&=1-P(\overline A)P(\overline B)\\
&=1-\frac12\cdot\frac56
=\frac7{12}.
\end{aligned}
\]
故选 C.
\end{solution}

\begin{problem}
已知 \(\triangle ABC\) 和点 \(M\) 满足 \(\overrightarrow{MA} + \overrightarrow{MB} + \overrightarrow{MC} = 0\). 若存在实数 \(m\) 使得 \(\overrightarrow{AB} + \overrightarrow{AC} = m\overrightarrow{AM}\) 成立，则 \(m =\)（\quad）

\choices
    {2}
    {3}
    {4}
    {5}
\end{problem}

\begin{answer}
B
\end{answer}

\begin{solution}
设点 \(A\)，\(B\)，\(C\)，\(M\) 的位置向量分别为
\(\bs{a}\)，\(\bs{b}\)，\(\bs{c}\)，\(\bs{m}\).由条件
\[
(\bs{a}-\bs{m})+(\bs{b}-\bs{m})
+(\bs{c}-\bs{m})=0
\]
得
\[
\bs{a}+\bs{b}+\bs{c}=3\bs{m}
\]
于是
\[
\begin{aligned}
\overrightarrow{AB}+\overrightarrow{AC}
&=(\bs{b}-\bs{a})+(\bs{c}-\bs{a})\\
&=\bs{b}+\bs{c}-2\bs{a}\\
&=3(\bs{m}-\bs{a})
=3\overrightarrow{AM}.
\end{aligned}
\]
与题设比较得到 \(m=3\)，故选 B.
\end{solution}

\begin{problem}
将参加夏令营的 600 名学生编号为：\(001\)、\(002\)，…，\(600\). 采用系统抽样方法抽取一个容量为 50 的样本，且随机抽得的号码为 \(003\). 这 600 名学生分住在三个营区，从 \(001\) 到 \(300\) 在第 \(I\) 营区，从 \(301\) 到 \(495\) 住在第 \(II\) 营区，从 \(496\) 到 \(600\) 住在第 \(III\) 营区. 三个营区被抽中的人数依次为（\quad）

\choices
    {26, 16, 8}
    {25, 17, 8}
    {25, 16, 9}
    {24, 17, 9}
\end{problem}

\begin{answer}
B
\end{answer}

\begin{solution}
系统抽样的抽样间隔为
\[
\frac{600}{50}=12
\]
随机抽到的第一个号码为 \(003\)，所以样本号码为
\(3\)，\(15\)，\(27\)，\(\ldots\)，\(3+49\times12=591\).
第 \(I\) 营区中的号码满足 \(3+12j\leqslant300\)，即
\(j=0\)，\(1\)，\(\ldots\)，\(24\)，共有 \(25\) 人.第 \(II\) 营区中的号码为
\(303\)，\(315\)，\(\ldots\)，\(495\)
共有 \(17\) 人；其余
\(507\)，\(519\)，\(\ldots\)，\(591\)
在第 \(III\) 营区，共有 \(8\) 人.因此人数依次为
\((25,17,8)\)，故选 B.
\end{solution}

\begin{problem}
如图，在半径为 \(r\) 的圆内作内接正六边形，再作正六边形的内切圆，又在此内切圆内作内接正六边形，如此无限继续下去，设 \( S_{n} \) 为前 \(n\) 个圆的面积之和，则 \( \lim\limits_{n\to\infty}S_{n}= \)（\quad）

\bitmapfigure[max width=0.38\linewidth]{img/2010/hubei_science/q7_fig1.png}

\choices
    {\(2\pi r^{2}\)}
    {\(\frac{8}{3}\pi r^2\)}
    {\(4\pi r^{2}\)}
    {\(6\pi r^{2}\)}
\end{problem}

\begin{answer}
C
\end{answer}

\begin{solution}
设第一个圆的半径为 \(r_1=r\)，以后各圆的半径依次为
\(r_2\)，\(r_3\)，\(\ldots\).正六边形的内切圆半径等于该正六边形外接圆半径的
\(\cos30^\circ\) 倍，所以
\(r_{n+1}=\frac{\sqrt3}{2}r_n\)，\(\pi r_{n+1}^2=\frac34\pi r_n^2\).
因此前 \(n\) 个圆的面积和是等比数列的前 \(n\) 项，首项为
\(\pi r^2\)，公比为 \(\frac34\).于是
\[
\lim_{n\to\infty}S_n
=\frac{\pi r^2}{1-\frac34}
=4\pi r^2
\]
故选 C.
\end{solution}

\begin{problem}
现安排甲，乙，丙，丁，戊5名同学参加上海世博会志愿者服务活动，每人从事翻译，导游，礼仪，司机四项工作之一，每项工作至少有一人参加. 甲，乙不会开车但能从事其他三项工作，丙，丁，戊都能胜任四项工作，则不同安排方案的种数是（\quad）

\choices
    {152}
    {126}
    {90}
    {54}
\end{problem}

\begin{answer}
B
\end{answer}

\begin{solution}
司机工作只能由丙、丁、戊承担.

若有两人从事司机工作，先从这三人中选出两人，有
\(\mathrm C_3^2\) 种；剩下三人分别从事另外三项工作，有 \(3!\) 种，共有
\[
\mathrm C_3^2\cdot3!=18
\]
种.

若只有一人从事司机工作，先选出司机，有 \(\mathrm C_3^1\) 种.剩下四人安排到翻译、导游、礼仪三项工作且每项至少一人，利用容斥原理，安排数为
\[
3^4-\mathrm C_3^1 2^4+\mathrm C_3^2 1^4
=81-48+3=36
\]
此时共有
\[
\mathrm C_3^1\cdot36=108
\]
种.因此总数为
\[
18+108=126
\]
故选 B.
\end{solution}

\begin{problem}
若直线 \( y = x + b \) 与曲线 \( y = 3 - \sqrt{4x - x^2} \) 有公共点，则 \( b \) 的取值范围是（\quad）

\choices
    {\([-1, 1 + 2\sqrt{2}]\)}
    {\([1 - 2\sqrt{2}, 1 + 2\sqrt{2}]\)}
    {\([1 - 2\sqrt{2},3]\)}
    {\([1 - \sqrt{2},3]\)}
\end{problem}

\begin{answer}
C
\end{answer}

\begin{solution}
曲线的定义域为 \(0\leqslant x\leqslant4\).直线与曲线有公共点时，公共点的横坐标 \(x\) 满足
\[
b=3-x-\sqrt{4x-x^2}
\]
令 \(t=x-2\)，则 \(t\in[-2,2]\)，上式右端化为
\[
\psi(t)=1-t-\sqrt{4-t^2}
\]
在 \((-2,2)\) 内，
\[
\psi'(t)=-1+\frac{t}{\sqrt{4-t^2}}
\]
唯一的驻点为 \(t=\sqrt2\)，且函数在此处由减变增，取得最小值
\[
\psi(\sqrt2)=1-2\sqrt2
\]
端点处
\(\psi(-2)=3\)，\(\psi(2)=-1\)
故最大值为 \(3\).由连续性，\(b\) 的取值范围为
\[
1-2\sqrt2\leqslant b\leqslant3
\]
故选 C.
\end{solution}

\begin{problem}
记实数 \(x_{1}\)，\(x_{2}\)，\(\cdots\)，\(x_{n}\) 中的最大数为 \(\max \{x_{1}, x_{2}, \cdots, x_{n}\}\)，最小数为 \(\min \{x_{1}, x_{2}, \cdots, x_{n}\}\). 已知 \(\triangle ABC\) 的三边边长为 \(a\)，\(b\)，\(c\)（\(a \leqslant b \leqslant c\)），定义它的倾斜度为 \(t = \max \left\{\frac{a}{b}, \frac{b}{c}, \frac{c}{a}\right\} \cdot \min \left\{\frac{a}{b}, \frac{b}{c}, \frac{c}{a}\right\}\)，则 \(t = 1\) 是 “\(\triangle ABC\) 为等边三角形”的（\quad）

\choices
    {必要而不充分的条件}
    {充分而不必要的条件}
    {充要条件}
    {既不充分也不必要的条件}
\end{problem}

\begin{answer}
A
\end{answer}

\begin{solution}
因为 \(a\leqslant b\leqslant c\)，所以
\(\frac ca\) 是三个比值中的最大值，而最小值是
\[
\min\left\{\frac ab,\frac bc\right\}
\]
若 \(a=b=c\)，则三个比值都为 \(1\)，从而 \(t=1\)，所以
\(t=1\) 是等边三角形的必要条件.

但它不是充分条件.例如取 \(a=b=2\)，\(c=3\)，这三数满足三角形不等式，构成等腰但非等边三角形，且
\[
\max\left\{1,\frac23,\frac32\right\}
\min\left\{1,\frac23,\frac32\right\}
=\frac32\cdot\frac23=1
\]
故 \(t=1\) 必要而不充分，选 A.
\end{solution}

\section{填空题}

\begin{problem}
在 \((x + \sqrt[4]{3} y)^{20}\) 的展开式中，系数为有理数的项共有\fillinblank{} 项.
\end{problem}

\begin{answer}
6
\end{answer}

\begin{solution}
展开式的通项为
\(T_{k+1}=\mathrm C_{20}^k x^{20-k}(\sqrt[4]{3}y)^k =\mathrm C_{20}^k3^{k/4}x^{20-k}y^k\)，\(k=0\)，\(1\)，\(\ldots\)，\(20\).
由于 \(\mathrm C_{20}^k\) 是整数，系数为有理数当且仅当
\(3^{k/4}\) 为有理数，即 \(k\) 是 \(4\) 的倍数.因此
\(k=0\)，\(4\)，\(8\)，\(12\)，\(16\)，\(20\)
共有 \(6\) 项.
\end{solution}

\begin{problem}
已知 \(z = 2x - y\)，式中变量 \(x\)，\(y\) 满足约束条件 \(\begin{cases}
y \leqslant x,\\
x + y \geqslant 1,\\
x \leqslant 2,
\end{cases}\) 则 \(z\) 的最大值为\fillinblank{}.
\end{problem}

\begin{answer}
5
\end{answer}

\begin{solution}
由 \(x+y\geqslant1\) 得 \(y\geqslant1-x\).对于固定的 \(x\)，
\(z=2x-y\) 随 \(y\) 增大而减小，所以要使 \(z\) 最大，应取最小的允许值
\(y=1-x\).此时
\[
z=2x-(1-x)=3x-1
\]
条件 \(y\leqslant x\) 与 \(y=1-x\) 给出 \(x\geqslant\frac12\)，再结合
\(x\leqslant2\)，可知 \(x\in[\frac12,2]\).故
\[
z=3x-1\leqslant3\cdot2-1=5
\]
当 \(x=2\)，\(y=-1\) 时，三个约束均满足，且 \(z=5\)，所以最大值为 \(5\).
\end{solution}

\begin{problem}
圆柱形容器内部盛有高度为 \(8 \mathrm{~cm}\) 的水，若放入三个相同的球 （球半径与圆柱的底面半径相同） 后，水恰好淹没最上面的球 （如图所示）. 则球的半径是\fillinblank{} \(\mathrm{cm}\).
\bitmapfigure[max width=0.24\linewidth]{img/2010/hubei_science/q13_fig1.png}
\end{problem}

\begin{answer}
4
\end{answer}

\begin{solution}
设球的半径以及圆柱的底面半径均为 \(r\).三个球竖直相切并放在容器底部，最上面球的顶部高度为 \(6r\).水恰好淹没最上面的球，所以放球后水和球共同占据的圆柱部分高度为 \(6r\)，三个球均完全浸没.

由体积守恒，放球前水的体积等于放球后圆柱体积减去三个球的体积，即
\[
8\pi r^2
=6\pi r^3-3\cdot\frac43\pi r^3
=2\pi r^3
\]
因 \(r>0\)，得 \(r=4\)，所以球的半径为 \(4\mathrm{~cm}\).
\end{solution}

\begin{problem}
某射手射击所得环数 \(\xi\) 的分布列如下：

\begin{center}
\begin{tabular}{c|cccc}
\(\xi\) & \(7\) & \(8\) & \(9\) & \(10\) \\
\hline
\(P\) & \(x\) & \(0.1\) & \(0.3\) & \(y\)
\end{tabular}
\end{center}
已知 \(\xi\) 的期望 \(E(\xi) = 8.9\)，则 \(y\) 的值为\fillinblank{}.
\end{problem}

\begin{answer}
0.4
\end{answer}

\begin{solution}
由概率和为 \(1\)，得
\(x+0.1+0.3+y=1\)，\(x+y=0.6\).
由期望公式，
\[
7x+8\times0.1+9\times0.3+10y=8.9
\]
即
\[
7x+10y=5.4
\]
与 \(x=0.6-y\) 联立，得
\(7(0.6-y)+10y=5.4\)，\(y=0.4\).
\end{solution}

\begin{problem}
设 \(a > 0\)，\(b > 0\). 称 \(\frac{2ab}{a + b}\) 为 \(a\)，\(b\) 的调和平均数. 如图，\(C\) 为线段 \(AB\) 上的点，且 \(AC = a\)，\(CB = b\)，\(O\) 为 \(AB\) 中点，以 \(AB\) 为直径做半圆. 过点 \(C\) 作 \(AB\) 的垂线交半圆于 \(D\). 连接 \(OD\)，\(AD\)，\(BD\). 过点 \(C\) 作 \(OD\) 的垂线，垂足为 \(E\). 则图中线段 \(OD\) 的长度是 \(a\)，\(b\) 的算术平均数，线段\fillinblank{} 的长度是 \(a\)，\(b\) 的几何平均数，线段\fillinblank{} 的长度是 \(a\)，\(b\) 的调和平均数.

\bitmapfigure[max width=0.42\linewidth]{img/2010/hubei_science/q15_fig1.png}
\end{problem}

\begin{answer}
\(CD\)；\(DE\)
\end{answer}

\begin{solution}
设 \(AB=a+b\)，取 \(A\) 为原点，令 \(AB\) 为 \(x\) 轴，则 \(C\) 的横坐标为
\(a\)，\(O\) 的横坐标为 \(\frac{a+b}{2}\).因为 \(CD\perp AB\)，且 \(D\) 在以
\(AB\) 为直径的半圆上，由圆的几何性质
\[
CD^2=AC\cdot CB=ab
\]
所以 \(CD=\sqrt{ab}\)，即为几何平均数.

又因为 \(OD\) 是半圆半径，
\[
OD=\frac{AB}{2}=\frac{a+b}{2}
\]
即为算术平均数.

在直角三角形 \(OCD\) 中，\(CE\perp OD\)，由射影定理
\[
OD\cdot DE=CD^2
\]
因此
\[
DE=\frac{CD^2}{OD}
=\frac{ab}{(a+b)/2}
=\frac{2ab}{a+b}
\]
即为调和平均数.故两空依次为 \(CD\)、\(DE\).
\end{solution}

\section{解答题}

\begin{problem}
已知函数 \(f(x) = \cos \left( \frac{\pi}{3} + x \right) \cos \left( \frac{\pi}{3} - x \right)\)，\(g(x) = \frac{1}{2} \sin 2x - \frac{1}{4}\).

\begin{enumerate}
\item 求函数 \( f(x) \) 的最小正周期；
\item 求函数 \( h(x) = f(x) - g(x) \) 的最大值，并求使 \( h(x) \) 取得最大值的 \(x\) 的集合.
\end{enumerate}
\end{problem}

\begin{solution}
\begin{enumerate}
\item 由积化和差公式，
\[
f(x)=\frac12\left[\cos\left(\frac{2\pi}{3}\right)+\cos2x\right]
=\frac12\cos2x-\frac14
\]
函数 \(\cos2x\) 的最小正周期为 \(\pi\)，乘以非零常数并作常数平移不改变周期，所以 \(f(x)\) 的最小正周期为 \(\pi\).

\item 由上式和 \(g(x)=\frac12\sin2x-\frac14\)，
\[
h(x)=f(x)-g(x)
=\frac12(\cos2x-\sin2x)
=\frac{\sqrt2}{2}\cos\left(2x+\frac\pi4\right)
\]
因此
\[
h(x)\leqslant\frac{\sqrt2}{2}
\]
最大值为 \(\frac{\sqrt2}{2}\).等号成立当且仅当
\(2x+\frac\pi4=2k\pi\)，\(k\in\mathbb Z\)
即
\(x=k\pi-\frac\pi8\)，\(k\in\mathbb Z\).
\end{enumerate}
\end{solution}

\begin{problem}
为了在夏季降温和冬季供暖时减少能源损耗，房屋的房顶和外墙需要建造隔热层. 某幢建筑物要建造可使用 \(20\text{ 年}\) 的隔热层，每厘米厚的隔热层建造成本为 \(6\text{ 万元}\). 该建筑物每年的能源消耗费用 \(C\)（单位：万元）与隔热层厚度 \(x\)（单位：cm）满足关系：\( C(x) = \frac{k}{3x + 5} \ (0 \leqslant x \leqslant 10) \)，若不建隔热层，每年能源消耗费用为 \(8\text{ 万元}\). 设 \( f(x) \) 为隔热层建造费用与 \(20\text{ 年}\)的能源消耗费用之和.

\begin{enumerate}
\item 求 \(k\) 的值及 \( f(x) \) 的表达式.
\item 隔热层修建多厚时，总费用 \( f(x) \) 达到最小，并求最小值.
\end{enumerate}
\end{problem}

\begin{solution}
\begin{enumerate}
\item 不建隔热层时 \(x=0\)，所以
\[
C(0)=\frac{k}{5}=8
\]
得 \(k=40\).厚度为 \(x\) cm 时，隔热层建造费用为 \(6x\) 万元，\(20\) 年能源消耗费用为
\[
20C(x)=\frac{800}{3x+5}
\]
故
\(f(x)=6x+\frac{800}{3x+5}\)，\(0\leqslant x\leqslant10\).

\item 对 \(f(x)\) 求导，得
\[
f'(x)=6-\frac{2400}{(3x+5)^2}
\]
在给定区间内 \(3x+5>0\)，所以
\[
f'(x)=0
\Longleftrightarrow (3x+5)^2=400
\Longleftrightarrow x=5
\]
当 \(0\leqslant x<5\) 时 \(3x+5<20\)，故 \(f'(x)<0\)；当
\(5<x\leqslant10\) 时 \(3x+5>20\)，故 \(f'(x)>0\).因此 \(f(x)\) 在
\([0,5]\) 上递减、在 \([5,10]\) 上递增，最小值在 \(x=5\) 处取得，且
\[
f(5)=6\times5+\frac{800}{20}=70
\]
所以隔热层厚 \(5\mathrm{~cm}\) 时总费用最小，最小值为 \(70\) 万元.
\end{enumerate}
\end{solution}

\begin{problem}
如图，在四面体 \(ABOC\) 中，\(OC \perp OA\)，\(OC \perp OB\)，\(\angle AOB = 120^{\circ}\)，且 \(OA = OB = OC = 1\).

\begin{enumerate}
\item 设 \(P\) 为 \(AC\) 的中点，证明：在 \(AB\) 上存在一点 \(Q\)，使 \(PQ \perp OA\)，并计算 \(\frac{AB}{AQ}\) 的值；
\item 求二面角 \(O - AC - B\) 的平面角的余弦值.
\end{enumerate}

\bitmapfigure[max width=0.38\linewidth]{img/2010/hubei_science/q18_fig1.png}
\end{problem}

\begin{solution}
建立空间直角坐标系，使 \(O\) 为原点，\(OA\) 在 \(x\) 轴正方向，\(OC\) 在
\(z\) 轴正方向.由已知可取
\(A(1,0,0)\)，\(B\left(-\frac12,\frac{\sqrt3}{2},0\right)\)，\(C(0,0,1)\).

\begin{enumerate}
\item 设 \(Q\) 将线段 \(AB\) 按比例 \(t\) 分割，即
\(Q=A+t\overrightarrow{AB} =\left(1-\frac32t,\frac{\sqrt3}{2}t,0\right)\)，\(0\leqslant t\leqslant1\).
点 \(P\) 为 \(AC\) 的中点，所以 \(P(\frac12,0,\frac12)\).于是
\(\overrightarrow{PQ} =\left(\frac12-\frac32t,\frac{\sqrt3}{2}t,-\frac12\right)\)，\(\overrightarrow{OA}=(1,0,0)\).
条件 \(PQ\perp OA\) 等价于
\[
\overrightarrow{PQ}\cdot\overrightarrow{OA}=0
\]
从而
\(\frac12-\frac32t=0\)，\(t=\frac13\).
这个 \(t\) 属于 \([0,1]\)，所以确有这样的点 \(Q\) 在 \(AB\) 上.又
\(AQ=t\,AB\)，故
\[
\frac{AB}{AQ}=\frac1t=3
\]

\item 平面 \(OAC\) 的一个法向量可取
\(\bs{n}_1=(0,1,0)\).平面 \(ABC\) 的两个方向向量为
\(\overrightarrow{AB}=\left(-\frac32,\frac{\sqrt3}{2},0\right)\)，\(\overrightarrow{AC}=(-1,0,1)\)
故其法向量可取
\(\bs{n}_2=\overrightarrow{AB}\times\overrightarrow{AC} =\left(\frac{\sqrt3}{2},\frac32,\frac{\sqrt3}{2}\right)\)，\(|\bs{n}_2|=\frac{\sqrt{15}}2\).
二面角的平面角取两平面的锐角，所以其余弦为两法向量夹角余弦的绝对值.于是
\[
\cos\theta
=\frac{|\bs{n}_1\cdot\bs{n}_2|}
{|\bs{n}_1|\,|\bs{n}_2|}
=\frac{3/2}{\sqrt{15}/2}
=\frac{\sqrt{15}}5
\]
\end{enumerate}
\end{solution}

\begin{problem}
已知一条曲线 \(C\) 在 \(y\) 轴的右边，\(C\) 上每一点到点 \(F(1,0)\) 的距离减去它到 \(y\) 轴距离的差都是 1.

\begin{enumerate}
\item 求曲线 \(C\) 的方程；
\item 是否存在正数 \(m\)，对于过点 \(M(m,0)\) 且与曲线 \(C\) 有两个交点 \(A\)，\(B\) 的任一直线，都有 \(\overrightarrow{FA} \cdot \overrightarrow{FB} < 0\)？若存在，求出 \(m\) 的取值范围；若不存在，请说明理由.
\end{enumerate}
\end{problem}

\begin{solution}
\begin{enumerate}
\item 设 \(C\) 上任一点为 \(P(x,y)\).因 \(P\) 在 \(y\) 轴右侧，\(x>0\)，且点
\(P\) 到 \(y\) 轴的距离为 \(x\).由题意，
\[
\sqrt{(x-1)^2+y^2}-x=1
\]
于是 \(\sqrt{(x-1)^2+y^2}=x+1>0\)，两边平方得
\((x-1)^2+y^2=(x+1)^2\)，\(y^2=4x\).
反过来，若 \(x>0\) 且 \(y^2=4x\)，则
\[
\sqrt{(x-1)^2+y^2}=\sqrt{(x+1)^2}=x+1
\]
满足原条件.因此曲线 \(C\) 的方程为
\(y^2=4x\)，\(x>0\).

\item 先设直线不垂直于 \(x\) 轴，写成 \(y=k(x-m)\).设它与 \(C\) 的两个交点为
\(A(x_1,y_1)\)、\(B(x_2,y_2)\)，则
\[
k^2(x-m)^2=4x
\]
即
\[
k^2x^2-(2mk^2+4)x+m^2k^2=0
\]
由韦达定理，
\(x_1+x_2=2m+\frac4{k^2}\)，\(x_1x_2=m^2\).
又 \(y_i=k(x_i-m)\)，所以
\[
y_1y_2=k^2(x_1-m)(x_2-m)=-4m
\]
因此
\[
\begin{aligned}
\overrightarrow{FA}\cdot\overrightarrow{FB}
&=(x_1-1)(x_2-1)+y_1y_2\\
&=x_1x_2-(x_1+x_2)+1+y_1y_2\\
&=m^2-6m+1-\frac4{k^2}.
\end{aligned}
\]
若 \(m^2-6m+1<0\)，上式必小于 \(0\)，所以所有这类直线都满足题设不等式.

还需考虑垂直于 \(x\) 轴的直线 \(x=m\).它与 \(C\) 的两个交点为
\((m,2\sqrt m)\)、\((m,-2\sqrt m)\)，且
\[
\overrightarrow{FA}\cdot\overrightarrow{FB}
=(m-1)^2-4m=m^2-6m+1
\]
因此题设的严格不等式对任意直线成立，当且仅当
\[
m^2-6m+1<0
\]
解得
\[
3-2\sqrt2<m<3+2\sqrt2
\]
由于 \(3-2\sqrt2>0\)，这正是所求正数 \(m\) 的取值范围.
\end{enumerate}
\end{solution}

\begin{problem}
已知数列 \(\{a_{n}\}\) 满足： \(a_1 = \frac{1}{2}\)，\(\frac{3(1 + a_{n + 1})}{1 - a_n} = \frac{2(1 + a_n)}{1 - a_{n + 1}}\)，\(a_n a_{n + 1} < 0\) （\(n \geqslant 1\)）；数列 \(\{b_n\}\) 满足： \(b_n = a_{n+1}^2 - a_n^2\) （\(n \geqslant 1\)）.

\begin{enumerate}
\item 求数列 \(\{a_{n}\}\)， \(\{b_{n}\}\) 的通项公式；
\item 证明：数列 \(\{b_n\}\) 中的任意三项不可能成等差数列.
\end{enumerate}
\end{problem}

\begin{solution}
\begin{enumerate}
\item 由题设递推式交叉相乘，得
\[
3(1+a_{n+1})(1-a_{n+1})
=2(1+a_n)(1-a_n)
\]
即
\[
1-a_{n+1}^2=\frac23(1-a_n^2)
\]
令 \(c_n=1-a_n^2\)，则 \(c_{n+1}=\frac23c_n\)，且
\(c_1=1-\frac14=\frac34\).所以
\(c_n=\frac34\left(\frac23\right)^{n-1}\)，\(a_n^2=1-\frac34\left(\frac23\right)^{n-1}\).
由 \(a_1>0\) 以及 \(a_na_{n+1}<0\)，数列 \(\{a_n\}\) 的符号交替，故
\[
a_n=(-1)^{n-1}
\sqrt{1-\frac34\left(\frac23\right)^{n-1}}
\]
再由 \(b_n=a_{n+1}^2-a_n^2\)，
\[
\begin{aligned}
b_n
&=\left[1-\frac34\left(\frac23\right)^n\right]
-\left[1-\frac34\left(\frac23\right)^{n-1}\right]\\
&=\frac14\left(\frac23\right)^{n-1}.
\end{aligned}
\]

\item 由 \(b_n=\frac14\left(\frac23\right)^{n-1}\)，数列 \(\{b_n\}\) 严格递减.
任取三个不同的项，设它们的下标为 \(i<j<k\)，并令
\(p=j-i\)、\(r=k-i\)，则 \(1\leqslant p<r\).若三项成等差数列，由大小关系知
中间项为 \(b_j\)，从而
\(2b_j=b_i+b_k\)，\(2\left(\frac23\right)^p=1+\left(\frac23\right)^r\).
右端大于 \(1\)，所以 \(2(\frac23)^p>1\)，只能有 \(p=1\).此时
\(\left(\frac23\right)^r=\frac13\)，\(2^r=3^{r-1}\)
但 \(r\geqslant2\) 时左边为偶数、右边为奇数，矛盾.因此
\(\{b_n\}\) 中任意三项都不可能成等差数列.
\end{enumerate}
\end{solution}

\begin{problem}
已知函数 \( f(x) = ax + \frac{b}{x} + c \)（\(a > 0\)）的图象在点\((1, f(1))\)处的切线方程为 \( y = x - 1 \).

\begin{enumerate}
\item 用 \(a\) 表示出 \(b\)，\(c\)；
\item 若 \( f(x) \geqslant \ln x \) 在 \([1, +\infty)\) 上恒成立，求 \(a\) 的取值范围.
\item 证明： \(1 + \frac{1}{2} + \frac{1}{3} + \cdots + \frac{1}{n} > \ln (n + 1) + \frac{n}{2(n + 1)}\)（\(n \geqslant 1\)）.
\end{enumerate}
\end{problem}

\begin{solution}
\begin{enumerate}
\item 因为切线 \(y=x-1\) 在 \(x=1\) 处，所以切点的纵坐标为 \(0\)，且切线斜率为 \(1\).由此
\(f(1)=a+b+c=0\)，\(f'(1)=a-b=1\).
解得
\(b=a-1\)，\(c=1-2a\).

\item 代入上面的 \(b\)，\(c\)，得
\[
f(x)=ax+\frac{a-1}{x}+1-2a
\]
令 \(G_a(x)=f(x)-\ln x\).将 \(a=\frac12\) 的情形分离出来，有
\[
G_a(x)=\left(\frac12x-\frac1{2x}-\ln x\right)
+\left(a-\frac12\right)\left(x+\frac1x-2\right)
\]
记
\[
G_{1/2}(x)=\frac12x-\frac1{2x}-\ln x
\]
则 \(G_{1/2}(1)=0\)，且
\[
G'_{1/2}(x)
=\frac12+\frac1{2x^2}-\frac1x
=\frac{(x-1)^2}{2x^2}\geqslant0
\]
所以当 \(x\geqslant1\) 时，\(G_{1/2}(x)\geqslant0\).又
\[
x+\frac1x-2=\frac{(x-1)^2}{x}\geqslant0
\]
故 \(a\geqslant\frac12\) 时 \(G_a(x)\geqslant0\)，即题设不等式恒成立.

反之，若 \(0<a<\frac12\)，则
\[
\lim_{x\to1^+}
\frac{G_{1/2}(x)}{x+\frac1x-2}
=\lim_{x\to1^+}
\frac{(x-1)^2/(2x^2)}{1-1/x^2}
=\lim_{x\to1^+}\frac{x-1}{2(x+1)}
=0
\]
因此当 \(x>1\) 充分接近 \(1\) 时，
\[
\frac{G_a(x)}{x+1/x-2}<0
\]
从而 \(G_a(x)<0\)，不可能在 \([1,+\infty)\) 上恒成立.故
\[
a\in\left[\frac12,+\infty\right)
\]

\item 对任意 \(k\geqslant1\)，由
\[
t-\ln(1+t)=\int_0^t\frac{u}{1+u}\,\mathrm du
\]
以及 \(0\leqslant u\leqslant t\)，有
\(t-\ln(1+t)> \int_0^t\frac{u}{1+t}\,\mathrm du =\frac{t^2}{2(1+t)}\)，\((t>0)\).
取 \(t=\frac1k\)，得到
\[
\frac1k-\ln\left(1+\frac1k\right)
>\frac1{2k(k+1)}
\]
对 \(k=1\)，\(2\)，\(\ldots\)，\(n\) 求和，并注意
\[
\sum_{k=1}^n\ln\left(1+\frac1k\right)
=\ln(n+1)
\]
便有
\[
\begin{aligned}
1+\frac12+\cdots+\frac1n-\ln(n+1)
&>\frac12\sum_{k=1}^n\left(\frac1k-\frac1{k+1}\right)\\
&=\frac12\left(1-\frac1{n+1}\right)
=\frac{n}{2(n+1)}.
\end{aligned}
\]
即得所证不等式.
\end{enumerate}
\end{solution}
